\documentclass[10pt,letterpaper]{article}
\usepackage[utf8]{inputenc}
\usepackage[intlimits]{amsmath}
\usepackage{amsfonts}
\usepackage{amssymb}
\usepackage{ragged2e}
\usepackage[letterpaper, margin=1in]{geometry}
\usepackage{graphicx}
\usepackage{cancel}
\usepackage{mathtools}
\usepackage{tabularx}
\usepackage{arydshln}
\usepackage{tensor}
\usepackage{array}
\usepackage{xcolor}
\usepackage[boxed]{algorithm}
\usepackage[noend]{algpseudocode}
\usepackage{listings}
\usepackage{textcomp}
% \usepackage[pdf,tmpdir,singlefile]{graphviz}
\usepackage{mathrsfs}
\usepackage{bbm}
\usepackage{tikz}
\usepackage{tikz-cd}
\usepackage{enumitem}
\usepackage{arydshln}
\usepackage{relsize}
\usepackage{multicol}
\usepackage{scalerel}

\usetikzlibrary{bayesnet}
\setlist{noitemsep}

%%%%%%%%%%%%%%%%%%%%%%%%%%%%%
% Formatting commands
%%%%%%%%%%%%%%%%%%%%%%%%%%%%%
\newcommand{\n}{\hfill\break}
\newcommand{\up}{\vspace{-\baselineskip}}
\newcommand{\hangblock}[2]{\par\noindent\settowidth{\hangindent}{\textbf{#1: }}\textbf{#1: }\!\!\!#2}
\newcommand{\lemma}[1]{\hangblock{Lemma}{#1}}
\newcommand{\defn}[1]{\hangblock{Defn}{#1}}
\newcommand{\thm}[1]{\hangblock{Thm}{#1}}
\newcommand{\cor}[1]{\hangblock{Cor}{#1}}
\newcommand{\prop}[1]{\hangblock{Prop}{#1}}
\newcommand{\ex}[1]{\hangblock{Ex}{#1}}
\newcommand{\exer}[1]{\hangblock{Exer}{#1}}
\newcommand{\fact}[1]{\hangblock{Fact}{#1}}
\newcommand{\remark}[1]{\hangblock{Remark}{#1}}
\newcommand{\proven}{\;$\square$\n}
\newcommand{\problem}[1]{\par\noindent{#1}\n}
\newcommand{\problempart}[2]{\par\noindent\indent{}\settowidth{\hangindent}{\textbf{(#1)} \indent{}}\textbf{(#1) }\!\!\!#2\n}
\newcommand{\ptxt}[1]{\textrm{\textnormal{#1}}}
\newcommand{\inlineeq}[1]{\centerline{$\displaystyle #1$}}
\newcommand{\pageline}{\up\par\noindent\rule{\textwidth}{0.1pt}}

%%%%%%%%%%%%%%%%%%%%%%%%%%%%%
% Math commands
%%%%%%%%%%%%%%%%%%%%%%%%%%%%%
% Set Theory
\newcommand{\card}[1]{\left|#1\right|}
\newcommand{\set}[1]{\left\{#1\right\}}
\newcommand{\setmid}{\;\middle|\;}
\newcommand{\ps}[1]{\mathcal{P}\left(#1\right)}
\newcommand{\pfinite}[1]{\mathcal{P}^{\ptxt{finite}}\left(#1\right)}
\newcommand{\naturals}{\mathbb{N}}
\newcommand{\N}{\naturals}
\newcommand{\integers}{\mathbb{Z}}
\newcommand{\Z}{\integers}
\newcommand{\rationals}{\mathbb{Q}}
\newcommand{\Q}{\rationals}
\newcommand{\reals}{\mathbb{R}}
\newcommand{\R}{\reals}
\newcommand{\complex}{\mathbb{C}}
\newcommand{\C}{\complex}
\newcommand{\halfPlane}{\mathbb{H}}
\let\H\relax
\newcommand{\H}{\halfPlane}
\newcommand{\comp}{^{\complement}}
\DeclareMathOperator{\Hom}{Hom}
\newcommand{\Ind}{\mathbbm{1}}
\newcommand{\cut}{\setminus}
\DeclareMathOperator{\elem}{elem}

% Graph Theory
\let\deg\relax
\DeclareMathOperator{\deg}{deg}
\newcommand{\degp}{\ptxt{deg}^{+}}
\newcommand{\degn}{\ptxt{deg}^{-}}
\newcommand{\precdot}{\mathrel{\ooalign{$\prec$\cr\hidewidth\hbox{$\cdot\mkern0.5mu$}\cr}}}
\newcommand{\succdot}{\mathrel{\ooalign{$\cdot\mkern0.5mu$\cr\hidewidth\hbox{$\succ$}\cr\phantom{$\succ$}}}}
\DeclareMathOperator{\cl}{cl}
\DeclareMathOperator{\affdim}{affdim}

% Probability
\newcommand{\parSymbol}{\P}
\newcommand{\Prob}{\mathbb{P}}
\renewcommand{\P}{\Prob}
\newcommand{\Avg}{\mathbb{E}}
\newcommand{\E}{\Avg}
\DeclareMathOperator{\Var}{Var}
\DeclareMathOperator{\cov}{cov}
\DeclareMathOperator{\Unif}{Unif}
\DeclareMathOperator{\Binom}{Binom}
\newcommand{\CI}{\mathrel{\text{\scalebox{1.07}{$\perp\mkern-10mu\perp$}}}}

% Standard Math
\newcommand{\inv}{^{-1}}
\newcommand{\abs}[1]{\left|#1\right|}
\newcommand{\ceil}[1]{\left\lceil{}#1\right\rceil{}}
\newcommand{\floor}[1]{\left\lfloor{}#1\right\rfloor{}}
\newcommand{\conj}[1]{\overline{#1}}
\newcommand{\of}{\circ}
\newcommand{\tri}{\triangle}
\newcommand{\inj}{\hookrightarrow}
\newcommand{\surj}{\twoheadrightarrow}
\newcommand{\ndiv}{\nmid}
\renewcommand{\epsilon}{\varepsilon}
\newcommand{\divides}{\mid}
\newcommand{\ndivides}{\nmid}
\DeclareMathOperator{\lcm}{lcm}
\DeclareMathOperator{\sgn}{sgn}
\newcommand{\map}[4]{\!\!\!\begin{array}[t]{rcl}#1 & \!\!\!\!\to & \!\!\!\!#2\\ #3 & \!\!\!\!\mapsto & \!\!\!\!#4\end{array}}
\newcommand{\bigsum}[2]{\smashoperator[lr]{\sum_{\scalebox{#1}{$#2$}}}}
\DeclareMathOperator{\gcf}{gcf}

% Linear Algebra
\newcommand{\Id}{\textrm{\textnormal{Id}}}
\newcommand{\im}{\textrm{\textnormal{im}}}
\newcommand{\norm}[1]{\abs{\abs{#1}}}
\newcommand{\tpose}{^{T}\!}
\newcommand{\iprod}[1]{\left<#1\right>}
\newcommand{\giprod}{\iprod{\;\,,\;}}
\DeclareMathOperator{\tr}{tr}
\DeclareMathOperator{\trace}{tr}
\newcommand{\chgBasMat}[3]{\!\!\tensor*[_{#1}]{\left[#2\right]}{_{#3}}}
\newcommand{\vecBas}[2]{\tensor*[]{\left[#1\right]}{_{#2}}}
\DeclareMathOperator{\GL}{GL}
\DeclareMathOperator{\Mat}{Mat}
\DeclareMathOperator{\vspan}{span}
\DeclareMathOperator{\rank}{rank}
\newcommand{\V}[1]{\vec{#1}}
\DeclareMathOperator{\proj}{proj}
\DeclareMathOperator{\compProj}{comp}
\DeclareMathOperator{\row}{row}
\newcommand{\smallPMatrix}[1]{\paren{\begin{smallmatrix}#1\end{smallmatrix}}}
\newcommand{\smallBMatrix}[1]{\brack{\begin{smallmatrix}#1\end{smallmatrix}}}
\newcommand{\pmat}[1]{\begin{pmatrix}#1\end{pmatrix}}
\newcommand{\bmat}[1]{\begin{bmatrix}#1\end{bmatrix}}

% Multilinear Algebra
\newcommand{\Lsym}{\L}
\let\L\relax
\DeclareMathOperator{\L}{\mathscr{L}}
\DeclareMathOperator{\A}{\mathcal{A}}
\DeclareMathOperator{\Alt}{Alt}
\DeclareMathOperator{\Sym}{Sym}
\newcommand{\ot}{\otimes}
\newcommand{\ox}{\otimes}
\DeclareMathOperator{\asc}{asc}
\DeclareMathOperator{\asSet}{set}
\DeclareMathOperator{\sort}{sort}
\DeclareMathOperator{\ringA}{\mathring{A}}

% Topology
\newcommand{\closure}[1]{\overline{#1}}
\newcommand{\uball}{\mathcal{U}}
\DeclareMathOperator{\Int}{Int}
\DeclareMathOperator{\Ext}{Ext}
\DeclareMathOperator{\Bd}{Bd}
\DeclareMathOperator{\rInt}{rInt}
\DeclareMathOperator{\ch}{ch}
\DeclareMathOperator{\ah}{ah}
\newcommand{\LargerTau}{\mathlarger{\mathlarger{\mathlarger{\mathlarger{\tau}}}}}
\newcommand{\Tau}{\mathcal{T}}

% Analysis
\DeclareMathOperator{\Graph}{Graph}
\DeclareMathOperator{\epi}{epi}
\DeclareMathOperator{\hypo}{hypo}
\DeclareMathOperator{\supp}{supp}
\newcommand{\lint}[2]{\underset{#1}{\overset{#2}{{\color{black}\underline{{\color{white}\overline{{\color{black}\int}}\color{black}}}}}}}
\newcommand{\uint}[2]{\underset{#1}{\overset{#2}{{\color{white}\underline{{\color{black}\overline{{\color{black}\int}}\color{black}}}}}}}
\newcommand{\alignint}[2]{\underset{#1}{\overset{#2}{{\color{white}\underline{{\color{white}\overline{{\color{black}\int}}\color{black}}}}}}}
\newcommand{\extint}{\ptxt{ext}\int}
\newcommand{\extalignint}[2]{\ptxt{ext}\alignint{#1}{#2}}
\newcommand{\conv}{\ast}
\newcommand{\pd}[2]{\frac{\partial{}#1}{\partial{}#2}}
\newcommand{\del}{\nabla}
\DeclareMathOperator{\grad}{grad}
\DeclareMathOperator{\curl}{curl}
\let\div\relax
\DeclareMathOperator{\div}{div}
\DeclareMathOperator{\vol}{vol}

% Complex Analysis
\let\Re\relax
\DeclareMathOperator{\Re}{Re}
\let\Im\relax
\DeclareMathOperator{\Im}{Im}
\DeclareMathOperator{\Res}{Res}

% Abstract Algebra
\DeclareMathOperator{\ord}{ord}
\newcommand{\generated}[1]{\left<#1\right>}
\newcommand{\cycle}[1]{\smallPMatrix{#1}}
\newcommand{\id}{\ptxt{id}}
\newcommand{\iso}{\cong}
\DeclareMathOperator{\Aut}{Aut}
\DeclareMathOperator{\SL}{SL}
\DeclareMathOperator{\op}{op}
\newcommand{\isom}[4]{\!\!\!\begin{array}[t]{rcl}#1 & \!\!\!\!\overset{\sim}{\to} & \!\!\!\!#2\\ #3 & \!\!\!\!\mapsto & \!\!\!\!#4\end{array}}
\newcommand{\F}{\mathbb{F}}
\newcommand{\acts}{\;\reflectbox{\rotatebox[origin=c]{-90}{$\circlearrowright$}}\;}
\newcommand{\disjunion}{\mathrel{\text{\scalebox{1.07}{$\perp\mkern-10mu\perp$}}}}
\DeclareMathOperator{\SO}{SO}
\DeclareMathOperator{\stab}{stab}
\DeclareMathOperator{\nullity}{nullity}
\DeclareMathOperator{\Perm}{Perm}
\DeclareMathOperator{\nsubgp}{\vartriangleleft}
\DeclareMathOperator{\notnsubgp}{\ntriangleleft}
\newcommand{\presentation}[2]{\left<#1\;\middle|\;#2\right>}
\DeclareMathOperator{\Char}{char}
\DeclareMathOperator{\fchar}{char}

% Convex Optimization
\newcommand{\sectionSymbol}{\S}
\let\S\relax
\newcommand{\S}{\mathbb{S}}
\DeclareMathOperator{\dist}{dist}
\DeclareMathOperator{\dom}{dom}
\DeclareMathOperator{\diag}{diag}
\DeclareMathOperator{\ones}{\mathbbm{1}}
\newcommand{\minimizeOver}[3]{\begin{array}{rl}\underset{#1}{\ptxt{minimize}} & #2\\ \ptxt{subject to} & #3\end{array}}
\newcommand{\maximizeOver}[3]{\begin{array}{rl}\underset{#1}{\ptxt{maximize}} & #2\\ \ptxt{subject to} & #3\end{array}}
\newcommand{\minimizationProblem}[2]{\minimizeOver{}{#1}{#2}}
\newcommand{\maximizationProblem}[2]{\maximizeOver{}{#1}{#2}}
\newcommand{\minimizeOverUnconstrained}[2]{\begin{array}{rl}\underset{#1}{\ptxt{minimize}} & #2\end{array}}
\newcommand{\maximizeOverUnconstrained}[2]{\begin{array}{rl}\underset{#1}{\ptxt{maximize}} & #2\end{array}}
\newcommand{\minimizationUnconstrained}[1]{\minimizeOverUnconstrained{}{#1}}
\newcommand{\maximizationUnconstrained}[1]{\maximizeOverUnconstrained{}{#1}}
\DeclareMathOperator{\argmin}{argmin}
\DeclareMathOperator{\argmax}{argmax}

% Proofs
\newcommand{\st}{s.t.}
\newcommand{\unique}{!}
\newcommand{\iffdef}{\overset{\ptxt{def}}{\Leftrightarrow}}
\newcommand{\eqdef}{\overset{\ptxt{def}}{=}}
\newcommand{\eqVertical}{\rotatebox[origin=c]{90}{=}}
\newcommand{\mapsfrom}{\mathrel{\reflectbox{\ensuremath{\mapsto}}}}
\newcommand{\mapsdown}{\rotatebox[origin=c]{-90}{$\mapsto$}\mkern2mu}
\newcommand{\mapsup}{\rotatebox[origin=c]{90}{$\mapsto$}\mkern2mu}
\newcommand{\from}{\!\mathrel{\reflectbox{\ensuremath{\to}}}}

% Brackets
\newcommand{\paren}[1]{\left(#1\right)}
\renewcommand{\brack}[1]{\left[#1\right]}
\renewcommand{\brace}[1]{\left\{#1\right\}}
\newcommand{\ang}[1]{\left<#1\right>}

% Algorithms
\algrenewcommand{\algorithmiccomment}[1]{\hskip 1em \texttt{// #1}}
\algrenewcommand\algorithmicrequire{\textbf{Input:}}
\algrenewcommand\algorithmicensure{\textbf{Output:}}
\newcommand{\algP}{\ptxt{\textbf{P}}}
\newcommand{\algNP}{\ptxt{\textbf{NP}}}
\newcommand{\algNPC}{\ptxt{\textbf{NP-Complete}}}
\newcommand{\algNPH}{\ptxt{\textbf{NP-Hard}}}
\newcommand{\algEXP}{\ptxt{\textbf{EXP}}}

%%%%%%%%%%%%%%%%%%%%%%%%%%%%%
% Other commands
%%%%%%%%%%%%%%%%%%%%%%%%%%%%%
\newcommand{\flag}[1]{\textbf{\textcolor{red}{#1}}}
\newcommand{\uSym}{\u}
\let\u\relax
\newcommand{\u}[1]{\underline{#1}}
\newcommand{\bSym}{\b}
\let\b\relax
\newcommand{\b}[1]{\textbf{#1}}
\newcommand{\iSym}{\i}
\let\i\relax
\newcommand{\i}[1]{\textit{#1}}

%%%%%%%%%%%%%%%%%%%%%%%%%%%%%%%%%%%%%%%
% Make l's curvy in math environments %
%%%%%%%%%%%%%%%%%%%%%%%%%%%%%%%%%%%%%%%
\mathcode`l="8000
\begingroup
\makeatletter
\lccode`\~=`\l
\DeclareMathSymbol{\lsb@l}{\mathalpha}{letters}{`l}
\lowercase{\gdef~{\ifnum\the\mathgroup=\m@ne \ell \else \lsb@l \fi}}%
\endgroup

%%%%%%%%%%%%%%%%%%%%%%%%%
% Fix \vdots and \ddots %
%%%%%%%%%%%%%%%%%%%%%%%%%
\usepackage{letltxmacro}
\LetLtxMacro\orgvdots\vdots
\LetLtxMacro\orgddots\ddots

\makeatletter
\DeclareRobustCommand\vdots{%
	\mathpalette\@vdots{}%
}
\newcommand*{\@vdots}[2]{%
	% #1: math style
	% #2: unused
	\sbox0{$#1\cdotp\cdotp\cdotp\m@th$}%
	\sbox2{$#1.\m@th$}%
	\vbox{%
		\dimen@=\wd0 %
		\advance\dimen@ -3\ht2 %
		\kern.5\dimen@
		% remove side bearings
		\dimen@=\wd2 %
		\advance\dimen@ -\ht2 %
		\dimen2=\wd0 %
		\advance\dimen2 -\dimen@
		\vbox to \dimen2{%
			\offinterlineskip
			\copy2 \vfill\copy2 \vfill\copy2 %
		}%
	}%
}
\DeclareRobustCommand\ddots{%
	\mathinner{%
		\mathpalette\@ddots{}%
		\mkern\thinmuskip
	}%
}
\newcommand*{\@ddots}[2]{%
	% #1: math style
	% #2: unused
	\sbox0{$#1\cdotp\cdotp\cdotp\m@th$}%
	\sbox2{$#1.\m@th$}%
	\vbox{%
		\dimen@=\wd0 %
		\advance\dimen@ -3\ht2 %
		\kern.5\dimen@
		% remove side bearings
		\dimen@=\wd2 %
		\advance\dimen@ -\ht2 %
		\dimen2=\wd0 %
		\advance\dimen2 -\dimen@
		\vbox to \dimen2{%
			\offinterlineskip
			\hbox{$#1\mathpunct{.}\m@th$}%
			\vfill
			\hbox{$#1\mathpunct{\kern\wd2}\mathpunct{.}\m@th$}%
			\vfill
			\hbox{$#1\mathpunct{\kern\wd2}\mathpunct{\kern\wd2}\mathpunct{.}\m@th$}%
		}%
	}%
}
\makeatother

\newcommand{\B}{
	\begin{tikzpicture}
	\filldraw [fill=red, draw=black] (0, 0) rectangle (0.37, 0.45);
	\draw [line width=0.5mm, white ] (0.1,0.08) -- (0.1,0.38);
	\draw[line width=0.5mm, white ] (0.1, 0.35) .. controls (0.2, 0.35) and (0.4, 0.2625) .. (0.1, 0.225);
	\draw[line width=0.5mm, white ] (0.1, 0.225) .. controls (0.2, 0.225) and (0.4, 0.1625) .. (0.1, 0.1);
	\end{tikzpicture}
}

\author{Professor Andrew Snowden\\ \small\i{Transcribed by Thomas Cohn}}
\title{Math 493 Lecture 23}
\date{11/25/19} % Can also use \today

\begin{document}
\maketitle
\setlength\RaggedRightParindent{\parindent}
\RaggedRight

\subsection*{Representation Theory}

\par\noindent
Recall from last time,\n

\defn{
	For $G$ a group, a \textbf{representation} of $G$ is a pair $(V,\rho)$, where $V$ is a vector space, and\break $\rho:G\to\GL(V)$ is a group homomorphism. In our case, we assume $V$ is finite dimensional, and specifically, a $K$-vector space (for some field $K$).\n
}

\par\noindent
Our problem is now understanding what representations of $G$ look like.\n

\defn{
	If $V,W$ are two representations of $G$. A \textbf{map of representations} (also known as a \textbf{$G$-map}) is a linear map $f:V\to{}W$ \st{} $f(gv)=gf(v)$, $\forall{}g\in{}G,v\in{}V$.\n
}

\defn{
	An \textbf{isomorphism of representations} is a bijective map of representations.\n
}

\par\noindent
The usual constructions of linear algebra apply to representations. Let $V,W$ be representations of $G$.
\begin{itemize}
	\item $V\oplus{}W$ is naturally a representation of $G$ -- $g(v,w)=(gv,gw)$.
	\item $\Hom_{G}(V,W)=\set{f:V\to{}W\mid{}f\;\ptxt{is a $G$-map}}$ is a vector space.
	\item $\Hom(V,W)=\set{f:V\to{}W\mid{}f\;\ptxt{is a linear map}}$ is a vector space.
\end{itemize}

\par\noindent
In fact, $\Hom(V,W)$ is naturally a representation. Given $g\in{}G$, $f\in\Hom(V,W)$, define $gf\in\Hom(V,W)$ by $(gf)(v)=gf(g\inv{}v)$.\n

\defn{
	$V^{G}=\set{v\in{}V\mid{}gv=v,\forall{}g\in{}G}$ is the \textbf{space of invariants}.\n
}

\par\noindent
Fact: $\Hom_{G}(V,W)=\Hom(V,W)^{G}$.\n
Reason: Say $f\in\Hom(V,W)$. Then $f$ is $G$-invariant (i.e. $f\in\Hom(V,W)^{G}$) iff $gf=f$, $\forall{}g\in{}G$
{\setlength{\abovedisplayskip}{0pt}
\begin{align*}
	& \Leftrightarrow{}gf(g\inv{}v)=f(v),\forall{}g\in{}G,v\in{}V\\
	& \Leftrightarrow{}gf(v)=f(gv),\forall{}g,v\\
	& \Leftrightarrow{}f\ptxt{ is a $G$-map}
\end{align*}}

\defn{
	$V^{*}=\Hom(V,K)$ (where $K$ is the one dimensional trivial representation over our field) is the \textbf{dual representation}.\n
}

\par\noindent
Pick bases for $V$ and $W$, with $\dim{}V=n$, $\dim{}W=m$. The action of $G$ on $V$ corresponds to $\rho:G\to\GL_{n}(K)$, and the action of $G$ on $W$ corresponds to $\sigma:G\to\GL_{m}(K)$. Then $V\oplus{}W$ has dimension $n+m$, and the matrix for $(\rho\oplus\sigma)(g)$ is
\up\[
	\begin{array}{c}
		\\
		\\
		n\\
		m
	\end{array}
	\begin{array}{c}
		\begin{array}{cc}
			\phantom{\rho(g)} & \phantom{\sigma(g)}\\
			n & m
		\end{array}\\
		\left[
		\begin{array}{c|c}
			\rho(g) & 0\\ \hline
			0 & \sigma(g)
		\end{array}
		\right]
	\end{array}
\]

\par\noindent
$\Hom(V,W)\cong{}M_{m,n}(K)$, the set of $m\times{}n$ matrices. For $\mathcal{A}\in{}M_{m,n}(K)$, $g\cdot{}A=\sigma(g)A\rho(g)\inv$.\n

\par\noindent
If $V,W\subseteq{}U$ are subrepresentations, then
\begin{itemize}
	\item $V\cap{}W$ is a subrepresentation.
	\item $V+W=\set{v+w\mid{}v\in{}V,w\in{}W}$ is a subrepresentation.
	\item $U/V$ is naturally a representation -- $g(u+V)=gu+V$.
	\item $V$ and $W$ are \textbf{complementary} if $U=V\oplus{}W$, which is true iff $V\cap{}W=\set{0}$ and $V+W=U$.
\end{itemize}

\defn{
	A representation $V$ is called \textbf{irreducible} if $V\ne{}0$ and its only subrepresentations are $0$ and $V$.\n
}

\par\noindent
Our approach to understanding representations is as follows. First, we will try to understand irreducible representations, and then we will understand how a general representation is build out of irreducible representations.\n

\ex{
	$G=D_{n}=\set{a,b\mid{}a^{2},b^{n},(ab)^{2}}$. We have a 2D representation $\rho:G\to\GL_{2}(\R)$, where $a\mapsto\smallPMatrix{0 & 1\\ 1 & 0}$ and $b\mapsto\smallPMatrix{\cos\frac{2\pi}{n} & -\sin\frac{2\pi}{n}\\ \sin\frac{2\pi}{n} & \cos\frac{2\pi}{n}}$.\n
	Fact: $\rho$ is an irreducible representation. Why? Because $\rho(b)$ doesn't have any real eigenvalues (all lines are rotated), so there isn't any line in $\R^{2}$ that's fixed by $G$.\n
}

\ex{
	$G=S_{n}$, $V=\C^{n}$, $\sigma{}e_{i}=e_{\sigma(i)}$.\n
	This is \textit{not} irreducible. Let $L=\vspan(v)$, where $v=e_{1}+\cdots+e_{n}=\smallBMatrix{1 & \cdots & 1}\tpose$. Then $\sigma{}v=v$, $\forall{}\sigma\in{}S_{n}$. Thus, $L$ is a subrepresentation.\n
	\n
	Does $L$ have a complementary subrepresentation?\n
	Let $T:\C\to\C^{n}$, where $f\mapsto{}e_{1}+\cdots+e_{n}=v$. ($f$ is the basis vector for $\C$.) This is a map of representations -- $T(\sigma{}f)=T(f)=v$, and $\sigma{}T(f)=\sigma{}v=v$.\n
	Now consider $S:\C^{n}\to\C$, where $e_{i}\mapsto{}1$. $\ker(S)=\set{\sum_{i=1}^{n}a_{i}e_{i}\;\middle|\;\sum_{i=1}^{n}a_{i}=0}$ is a subrepresentation of $\C^{n}$, and has basis $(e_{1}-e_{n},e_{2}-e_{n},\ldots,e_{n-1}-e_{n})$.\n
	These are complementary subrepresentations, so $\C^{n}=\ker(S)\oplus{}L$.\n
}

\par\noindent
Fact: Not every subrepresentation has a complementary subrepresentation.\n

\ex{
	$\rho:\Z=G\to\GL_{2}$, where $n\mapsto\smallPMatrix{1 & n\\ 0 & 1}$. Note that $\smallPMatrix{1 & n\\ 0 & 1}\smallPMatrix{1 & m\\ 0 & 1}=\smallPMatrix{1 & n+m\\ 0 & 1}$. $L=\vspan(e_{1})=\smallBMatrix{*\\ 0}$ is a subrepresentation, but no other $1$D space is mapped to itself, so $L$ has no complementary subrepresentation.\n
}

\ex{
	$G=\Z/p\Z$ (for $p$ prime), $K=\F_{p}$.\n
	Let $\rho:G\to\GL_{2}(\F{n})$ where $n\mapsto\smallPMatrix{1 & n\\ 0 & 1}$. This is also an example, as $\vspan(e_{1})$ has no complement.\n
}

\thm{
	If $G$ is finite and $\fchar(K)=0$, then every subrepresentation has a complement.\n
	Proof: Let $V$ be some representation, and $W$ a subrepresentation of $V$. Let $\pi:V\to{}V/W$ be the projection map. Let $S:V/W\to{}V$ \st{} $\pi(S(x))=x$. Note that $S$ is not necessarily a $G$-map.\n
	We will show that, for any choice of $S$, $\im(S)$ is a complementary subspace to $W$: Given some $v\in{}V$, we have
	\[
		\pi(v-S(\pi(v)))=\pi(v)-\pi(S(\pi(v)))=\pi(v)-\pi(v)=0
	\]
	So $v-S(\pi(v))\in{}W=\ker(\pi)$. So $v=(\underbrace{v-S(\pi(v))}_{\in{}W})+\underbrace{S(\pi(v))}_{\in\im(S)}$.\n
	If $S$ is a $G$-map, then $\im(S)$ is a subrepresentation, and it's a complementary subrepresentation to $W$. To ensure $S$ is a $G$-map, we can average it. Specifically, define
	\[
		S':V/W\to{}V\qquad{}S'(x)=\frac{1}{\#G}\sum_{g\in{}G}gS(g\inv{}x)
	\]
	We claim that (1) $S'$ is a $G$-map, and (2), that $\pi\of{}S'=\Id$.
	\begin{enumerate}[topsep=0pt,leftmargin=4\parindent,label=(\arabic*)]
		\item For $h\in{}G$,
		\[
			S'(hx)=\frac{1}{\#G}\sum_{\substack{g\in{}G\\ g=hg'}}gS(g\inv{}hx)=\frac{1}{\#G}\sum_{g'\in{}G}(hg')S((g')\inv{}h\inv{}hx)=hg'(x)
		\]
		\item Because $\pi$ is a $G$-map,
		\[
			\pi(S'(x))=\pi\paren{\frac{1}{\#G}\sum_{g\in{}G}gS(g\inv{}x)}=\frac{1}{\#G}\sum_{g\in{}G}g\pi(S(g\inv{}x))=\frac{1}{\#G}\sum_{g\in{}G}g(g\inv(x))=x
		\]
	\end{enumerate}\up\n
	Thus, $\im(S')$ is a complementary subrepresentation to $W$.\proven
}

\par\noindent
Note: The theorem and proof are also valid if $p=\fchar(K)>0$ and $p\nmid\#{}G$.\n

\cor{
	Let $G$ be a finite group, $\fchar(K)=0$ (or it doesn't divide $G$). Then every finite dimensional representation of $G$ is the direct product of irreducible representations.\n
	Proof: We'll proceed by induction on the dimension of $V$. Let $V$ be given. If $V$ is irreducible, then there is nothing to prove. If $V=\set{0}$, then there is nothing to prove.\n
	Assume $V\ne\set{0}$, and $V$ is not irreducible. Then $\exists{}W\subset{}V$ a subrepresentation, with $W\ne\set{0},W\ne{}V$. By the above theorem, $\exists{}U$ a complementary subspace, so $V=U\oplus{}W$. By induction, $U$ and $W$ are the direct products of irreducible representations, so $V$ is too.\proven
}

\par\noindent
The previous remark is still true!\n

\par\noindent
Given the corollary, the main problem (for $G$ finite, $\fchar(K)=0$) is to understand the irreducible representations. There is a beautiful solution when $K=\C$ using \textit{character theory}. For now, assume $K=\C$, $\#G<\infty$.\n

\defn{
	Given a representation $V$ of $G$, we define its \textbf{character} to be the map $\map{\chi_{V}:G}{\C}{g}{\tr(g|_{V})}$\n ($g|_{V}$ is the linear operator $g$ acting on $V$.)\n
}

\par\noindent
We have the really cool fact that $V\cong{}W$ iff $\chi_{V}=\chi_{W}$.

\end{document}